\chapter{Wheeled Mobile Robots: Mechanics}
\label{ch:wmr-mechanics}

This chapter introduces the mechanical structure of wheeled mobile
robots. It will focus on what kinds of wheels exist, how they
can be assembled into kinematic structures, and what constraints are imposed on the motion of the robot
in various situations.

\section{Sources of Locomotion}
\label{sec:locomotion}

The main sources of movement for a mobile robot are:
\begin{itemize}
    \item \textbf{Wheels}: wheeled robots usually consist of a single
          rigid body, with the wheels in contact with the ground.
    \item \textbf{Feet}: legged robots, consisting of several rigid
          bodies articulated through joints.
\end{itemize}

In this course we focus on wheeled mobile robots (WMRs), which are the
simplest and most widely deployed class of mobile robots.

\section{Static Balance}
\label{sec:static-balance}

The wheels are in contact with the ground, and each wheel defines a
point of contact. The \textbf{support polygon} is the convex hull of
all the contact points.

\begin{definition}[Static Balance]
A wheeled robot is in \textbf{static balance} if the projection on the
ground of its center of mass falls inside the support polygon.
\end{definition}

A robot must be balanced in order not to fall. The convex hull of two
points is a line segment, hence a robot with only two wheels cannot be
easily balanced in a static fashion: it needs at least a third contact
point (typically a caster wheel). In the same way, a vehicle with a
single wheel (the unicycle) has a support polygon that degenerates to a
point, and a bicycle has a support polygon that degenerates to a line
segment: neither is statically balanced.

\begin{figure}[H]
    \centering
    \includegraphics[width=0.5\linewidth]{images/balance.pdf}
    \caption{Support polygon and static balance.}
    \label{fig:support-polygon}
\end{figure}

\section{Types of Wheels}
\label{sec:types-of-wheels}

There are three main types of wheels:

\begin{itemize}
    \item \textbf{Fixed wheels}: they have a fixed orientation with
          respect to the robot body. They can roll but cannot steer
          about their vertical axis. They are usually actuated; if
          passive, they are used for balance and offer a contact point.
    \item \textbf{Steerable wheels}: they can steer and are used for
          orientation, hence they are typically active (e.g.\ the front
          wheels of a car).
    \item \textbf{Caster wheels}: they have a variable orientation but
          are typically not actuated. Their axis of orientation is not
          aligned with the wheel axis, so they align automatically with
          the direction of motion. They are kinematically
          ``invisible'' but offer a contact point for balance.
\end{itemize}

\begin{figure}[ht]
    \centering
    \includegraphics[width=0.65\linewidth]{images/wheels.png}
    \caption{Fixed, steerable, and caster wheels.}
    \label{fig:wheel-types}
\end{figure}

\section{Kinematic Structures}
\label{sec:kinematic-structures}

\subsection{Differential Drive}
\label{subsec:differential-drive}

The simplest wheeled robot is the \textbf{differential drive}: two
aligned fixed wheels, independently actuated, plus one caster wheel for
balance. The two actuated wheels are driven independently, which
allows the robot to rotate in place by making them turn in opposite
directions.

\subsection{Synchro-Drive}
\label{subsec:synchro-drive}

A \textbf{synchro-drive} robot has three steerable wheels and two
actuators: one controls simultaneously the orientation of all the
wheels, and the other rolls them together. The base over the wheels
never changes orientation. This configuration is preferable when the
body to be moved is particularly heavy.

\begin{figure}[H]
    \centering
    \includegraphics[width=\linewidth]{images/structures.eps}
    \caption{Differential drive (left) and synchro-drive (right).}
    \label{fig:structures}
\end{figure}

\subsection{Tricycle and Car-Like}
\label{subsec:tricycle-car}

The \textbf{tricycle} has two fixed wheels at the back and one actuated
steerable wheel at the front. The \textbf{car-like} robot is similar,
but with two steerable wheels at the front.

Both are subject to \emph{nonholonomic} constraints: the pure rolling
condition on each wheel forbids sideways motion, but does not prevent
the robot from reaching any configuration in the plane.

\section{The Differential}
\label{sec:differential}

Consider a car-like robot taking a curve. In the same time interval,
the outer wheel travels a longer distance than the inner wheel, so its
rotation speed must be higher. If the two front wheels were rigidly
connected to a single motor, they would be forced to spin at the same
speed, causing tire slip and mechanical strain.

To avoid this, a device called \textbf{differential} is used.

\begin{definition}[Differential]
A mechanical differential is a gear assembly (typically bevel
or planetary gears) connected to a single drive motor. It splits the
power to the two wheels on the same axle while allowing them to rotate
at different speeds. When the robot turns, the outer wheel travels a
longer arc than the inner one; the differential lets each wheel spin
at its natural speed, preventing tire slip, scrubbing, and mechanical
strain on the drivetrain.
\end{definition}

\begin{figure}[ht]
    \centering
    \includegraphics[width=0.5\linewidth]{images/differential.eps}
    \caption{Path of a car-like robot turning: the outer wheel travels
    a longer distance than the inner one.}
    \label{fig:differential}
\end{figure}

\section{Constraints}
\label{sec:constraints}

The presence of wheels in contact with the ground introduces
\emph{constraints} on the motion of the robot. Two classes matter for
this course: geometric constraints, which act on the configuration, and
kinematic constraints, which act on the velocities.

\subsection{Geometric Constraints}
\label{subsec:geometric-constraints}

We focus on equality constraints that are time-invariant. A geometric
constraint has the form
\begin{equation}
    h(q) = 0,
\end{equation}
and depends only on the configuration. In general we have \( k < n \)
independent constraints,
\begin{align}
    h_1(q) &= 0, \\
           &\vdots \\
    h_k(q) &= 0.
\end{align}
If the configuration space is an \( n \)-manifold, \( k \) independent
geometric constraints reduce the accessible configurations to an
\( (n - k) \)-manifold embedded in \( C \), called a
\textbf{sub-manifold}.

For example, if \( C = \mathbb{R}^2 \) (a 2-manifold), one constraint
restricts the configurations to a curve in the plane. If
\( C = \mathbb{R}^3 \), one constraint restricts them to a surface.

\begin{nb}[Solving for a Coordinate]
Instead of using the full vector \( q \), one can solve
\( h(q) = 0 \) for one coordinate, e.g.\
\( q_i = \tilde{h}(q_1, \dots, q_{i-1}, q_{i+1}, \dots, q_n) \). But
this choice is not always good: for the constraint
\( q_1^2 + q_2^2 = 1 \), solving for \( q_2 \) gives
\( q_2 = \pm\sqrt{1 - q_1^2} \), which is not single-valued. For a
circle, it is better to use the angle \( \alpha \in \mathcal{S}^1 \)
as the single coordinate, with \( q_1 = \cos\alpha \) and
\( q_2 = \sin\alpha \).
\end{nb}

\subsection{Kinematic Constraints}
\label{subsec:kinematic-constraints}

A kinematic constraint is local: it prevents the configuration from
moving in certain directions, but does not necessarily restrict the
accessible configurations. A general kinematic constraint has the form
\begin{equation}
    a(q, \dot{q}) = 0,
\end{equation}
and depends on the state and its velocity. All the relevant kinematic
constraints in robotics are in \textbf{Pfaffian form}:
\begin{equation}
\label{eq:pfaffian}
    a^\top(q) \dot{q} = 0,
\end{equation}
with \( a(q) \) an \( n \)-dimensional vector depending on the
configuration. The Pfaffian form expresses an orthogonality condition:
at each configuration, the velocity must be orthogonal to \( a(q) \).

Multiple kinematic constraints can be arranged in matrix form:
\begin{equation}
\begin{pmatrix}
    a_1^\top(q) \\
    a_2^\top(q) \\
    \vdots \\
    a_k^\top(q)
\end{pmatrix}
\dot{q} = 0
\implies
A(q) \dot{q} = 0
\iff
\dot{q} \in \mathcal{N}(A(q)),
\end{equation}
where \( A(q) \) is a \( k \times n \) matrix. The admissible
velocities are the elements of the null space of \( A(q) \), a linear
subspace of \( \mathbb{R}^n \).

\begin{nb}[Geometric Implies Kinematic]
Any geometric constraint \( h(q) = 0 \) implies a kinematic
constraint. Differentiating with respect to time:
\[
\dot{h}(q) = \frac{\partial h}{\partial q} \dot{q}
           = (\nabla_q h)^\top \dot{q} = 0.
\]
This is exactly a Pfaffian constraint with
\( a(q) = \nabla_q h \). The gradient is orthogonal to the
sub-manifold defined by the geometric constraint, so the kinematic
constraint prevents the robot from leaving the sub-manifold.
\end{nb}

The converse is not generally true: a kinematic constraint
\( a^\top(q) \dot{q} = 0 \) implies a geometric constraint only if
the constraint is \emph{integrable}, i.e.\ only if there exist a
function \( h(q) \) and a nonvanishing function \( \lambda(q) \) such
that
\begin{equation}
\label{eq:integrability}
    \nabla_q h = \lambda(q)\, a(q).
\end{equation}
(The factor \( \lambda \) is needed because multiplying a constraint by
a nonzero function does not change the set of admissible velocities;
in the simplest case \( \lambda = 1 \).)

\subsection{Holonomic vs Nonholonomic Constraints}
\label{subsec:holonomic-nonholonomic}

\begin{definition}[Holonomic Constraint]
\label{def:holonomic}
An integrable kinematic constraint is called \textbf{holonomic}. It
defines a geometric constraint, and therefore reduces the dimension of
the accessible configuration space.
\end{definition}

\begin{definition}[Nonholonomic Constraint]
\label{def:nonholonomic}
A kinematic constraint that is not integrable is called
\textbf{nonholonomic}. It restricts the set of admissible velocities at
each configuration, without being equivalent to a geometric constraint.
Depending on the degree of integrability, the accessible configuration
space is either the whole \( C \) (\emph{completely nonholonomic}
constraints) or a sub-manifold of dimension strictly larger than the
number of independent admissible velocity directions (\emph{partially
integrable} constraints, see \cref{subsec:controllability}).
\end{definition}

\paragraph{Example: integrable constraint.}
Consider the kinematic constraint for a 2-dimensional robot:
\[
q_1 \dot{q}_1 + q_2 \dot{q}_2 = 0.
\]
This is integrable: the resulting geometric constraint is
\( q_1^2 + q_2^2 = c \), with \( c > 0 \). So it is holonomic.

\paragraph{Example: nonholonomic constraint.}
The simplest nonholonomic system is the \textbf{rolling disk}. Its
configuration is \( q = (x, y, \theta)^\top \), with \( \theta \) the
orientation with respect to the horizontal axis. In pure rolling (no
slippage), the disk can only move in the direction it points to, not
laterally. The constraint is
\begin{equation}
\label{eq:rolling}
    \sin\theta \, \dot{x} - \cos\theta \, \dot{y} = 0,
\end{equation}
or in vector form
\[
\begin{pmatrix} \sin\theta & -\cos\theta & 0 \end{pmatrix}
\begin{pmatrix} \dot{x} \\ \dot{y} \\ \dot{\theta} \end{pmatrix} = 0.
\]
The vector \( n^\top = (\sin\theta, -\cos\theta) \) is normal to the
direction of motion. Any velocity with a component along \( n \) is
not allowed. The line through the contact point with direction \( n \)
is called the \textbf{zero-motion line}: the velocity of the contact
point has zero component along it.

This constraint does not restrict the reachable configurations: every
configuration can still be reached. To see this, construct an explicit
trajectory (a \textbf{universal maneuver}) between any two
configurations \( q_s \) and \( q_g \):
\begin{enumerate}
    \item Rotate in place from \( \theta_s \) to the angle
          \( \theta^* \) such that the vector
          \( (x_g - x_s, y_g - y_s) \) is aligned with the sagittal
          axis.
    \item Move along that direction until \( (x_g, y_g) \) is reached.
    \item Rotate in place from \( \theta^* \) to \( \theta_g \).
\end{enumerate}
Each step respects the rolling constraint. So the rolling disk is
nonholonomic but controllable.

\begin{nb}[Holonomic vs Nonholonomic: The Takeaway]
\begin{itemize}
    \item \textbf{Holonomic} (integrable): the accessible
          configuration space is reduced to a sub-manifold. The robot
          is confined to it, hence the system is \emph{not}
          controllable.
    \item \textbf{Completely nonholonomic}: the accessible
          configuration space is the whole \( C \). The robot can
          reach any configuration, but only along admissible paths:
          the system is controllable.
    \item \textbf{Partially integrable} (still nonholonomic): the
          accessible configuration space is a sub-manifold, but larger
          than the one a holonomic constraint would give. The system
          is not controllable.
\end{itemize}
Controllability of the kinematic model is therefore equivalent to the
constraints being completely nonholonomic. The precise test is given in
\cref{subsec:controllability}.
\end{nb}

\section{Kinematic Models}
\label{sec:kinematic-models}

The kinematic constraints are written in Pfaffian form
\( A(q) \dot{q} = 0 \), with \( A(q) \) a \( k \times n \) matrix.
Admissible velocities belong to the null space of \( A(q) \). We want
to find a basis for this null space:
\[
\mathcal{N}(A(q)) = \text{span}\{ g_1(q), g_2(q), \dots, g_m(q) \},
\]
where \( m := n - k \). Each \( g_i(q) \) is a vector field on \( C \)
with values in \( \mathbb{R}^n \), and every admissible velocity is a
linear combination of the basis vectors:
\begin{equation}
    \dot{q} = \sum_{i=1}^{m} u_i \, g_i(q).
\end{equation}
In compact form, introducing the \( m \)-dimensional vector
\( u = (u_1, \dots, u_m)^\top \):
\begin{equation}
\label{eq:kinematic-model}
    \dot{q} = \underbrace{
    \begin{pmatrix}
        | & | & & | \\
        g_1(q) & g_2(q) & \dots & g_m(q) \\
        | & | & & |
    \end{pmatrix}}_{G(q)}
    \underbrace{
    \begin{pmatrix}
        u_1 \\ \vdots \\ u_m
    \end{pmatrix}}_{u}.
\end{equation}
This is the \textbf{kinematic model} of the robot, with \( u \) a
velocity-level quantity.

\paragraph{Driftless property.}
The kinematic model \cref{eq:kinematic-model} satisfies
\[
u = 0 \implies \dot{q} = 0.
\]
A system with this property is called \textbf{driftless}: it has no
free evolution. When the inputs are zero, the state does not move.

\paragraph{Pseudovelocities.}
The vector \( u \) does not coincide with the actual control inputs of
the robot (which are forces and torques), nor with the generalized
velocities \( \dot{q} \). It is a velocity-level quantity filtered
through the kinematic constraints, and it is called
\textbf{pseudo-velocity}. The kinematic model describes the motion as
resulting from the pseudo-velocities \( u \) filtered through the
constraints \( A(q) \).

\paragraph{Non-uniqueness of the model.}
The choice of the basis \( g_1, \dots, g_m \) is not unique, and
different choices lead to equivalent models with different physical
interpretations of \( u \). In practice, we choose the basis so that
the components of \( u \) have a natural meaning (e.g.\ linear and
angular velocity for a unicycle).

\subsection{Controllability and the Accessibility Rank Condition}
\label{subsec:controllability}

Whether the constraints are holonomic or nonholonomic can be decided by
studying the controllability of the kinematic model
\cref{eq:kinematic-model}. Two cases are possible:
\begin{itemize}
    \item the system is \textbf{controllable}: for any two
          configurations \( q_i, q_f \in C \) there exists an input
          \( u(t) \) that steers the system from \( q_i \) to
          \( q_f \) while satisfying the constraints. The constraints
          do not affect the accessibility of \( C \): they are
          (completely) nonholonomic;
    \item the system is \textbf{not controllable}: the constraints
          reduce the set of accessible configurations, so they are at
          least partially integrable.
\end{itemize}

To test controllability we need the Lie bracket of two vector fields,
\begin{equation}
\label{eq:lie-bracket}
    [g_i, g_j](q) = \frac{\partial g_j}{\partial q}\, g_i(q)
                  - \frac{\partial g_i}{\partial q}\, g_j(q),
\end{equation}
which is a new vector field, in general not a linear combination of
\( g_i \) and \( g_j \). Let \( \Delta = \text{span}\{g_1, \dots, g_m\} \)
be the distribution generated by the input vector fields. A
distribution is \textbf{involutive} if it is closed under Lie bracket.
The \textbf{accessibility distribution} \( \Delta_{\mathcal{A}} \) is
the involutive closure of \( \Delta \), i.e.\ the smallest involutive
distribution containing it. It is built iteratively,
\[
\Delta_1 = \Delta, \qquad
\Delta_{i+1} = \Delta_i + \text{span}\{ [g, h] \;:\; g \in \Delta_1,\ h \in \Delta_i \},
\]
until the dimension stops growing (\( \Delta_{\mathcal{A}} = \Delta_{i}
\) when \( \Delta_{i+1} = \Delta_i \)).

The controllability test is the \textbf{accessibility rank condition}:
\begin{equation}
\label{eq:arc}
    \dim \Delta_{\mathcal{A}}(q) = n.
\end{equation}
Let \( \nu = \dim \Delta_{\mathcal{A}} \) and \( m = \dim \Delta \).
Three cases:
\begin{itemize}
    \item \( \nu = n \): \cref{eq:arc} holds, the system is
          controllable, the constraints are \emph{completely
          nonholonomic};
    \item \( m < \nu < n \): the system is not controllable, but the
          loss of accessibility is not maximal; the constraints are
          \emph{partially integrable} (the mechanical system is still
          nonholonomic);
    \item \( \nu = m \): the system is not controllable and the loss of
          accessibility is maximal; the constraints are \emph{completely
          integrable}, hence holonomic. This happens exactly when
          \( \Delta_{\mathcal{A}} = \Delta \), i.e.\ when \( \Delta \)
          is already involutive.
\end{itemize}

The \textbf{degree of nonholonomy} \( \kappa \) of a controllable
system is the smallest integer such that \( \dim \Delta_\kappa = n \),
i.e.\ the minimum number of vector fields appearing in the nested
brackets needed to generate \( \mathbb{R}^n \). It quantifies \emph{how}
nonholonomic the constraints are.

\begin{nb}[Systems with a Single Input Vector Field]
If the number of Pfaffian constraints is \( k = n - 1 \), then
\( m = 1 \): the kinematic model has a single vector field, and
\( \Delta \) is always involutive (the bracket of a field with itself
is zero). Hence a set of \( n-1 \) Pfaffian constraints is always
completely integrable. For example, a scalar Pfaffian constraint is
always holonomic for a mechanical system with two generalized
coordinates.
\end{nb}

\subsection{The Unicycle Model}
\label{subsec:unicycle-model}

Consider the rolling disk from
\cref{subsec:holonomic-nonholonomic}. Its constraint is
\cref{eq:rolling}:
\[
\underbrace{
\begin{pmatrix} \sin\theta & -\cos\theta & 0 \end{pmatrix}
}_{A(q)}
\begin{pmatrix} \dot{x} \\ \dot{y} \\ \dot{\theta} \end{pmatrix} = 0.
\]
The dimension of \( \mathcal{N}(A) \) is \( m = n - k = 3 - 1 = 2 \).
A basis is
\[
g_1(q) = \begin{pmatrix} \cos\theta \\ \sin\theta \\ 0 \end{pmatrix},
\qquad
g_2(q) = \begin{pmatrix} 0 \\ 0 \\ 1 \end{pmatrix}.
\]
The kinematic model is
\begin{equation}
    \dot{q} = G(q) u =
    \begin{pmatrix}
        \cos\theta & 0 \\
        \sin\theta & 0 \\
        0 & 1
    \end{pmatrix}
    \begin{pmatrix} u_1 \\ u_2 \end{pmatrix}
    \implies
    \begin{cases}
        \dot{x} = u_1 \cos\theta, \\
        \dot{y} = u_1 \sin\theta, \\
        \dot{\theta} = u_2.
    \end{cases}
\end{equation}
The pseudo-velocity \( u_1 \) is the scalar velocity along the sagittal
axis:
\[
u_1 = \pm\sqrt{\dot{x}^2 + \dot{y}^2},
\]
while \( u_2 = \dot{\theta} \). Rename them:
\[
v = u_1, \qquad \omega = u_2.
\]
The linear velocity \( v \) is related to the rolling velocity
\( \omega_{\text{roll}} \) and the wheel radius by
\( v = \omega_{\text{roll}} \cdot r \). The model becomes
\begin{equation}
\label{eq:unicycle}
    \dot{q} = \begin{pmatrix} \cos\theta \\ \sin\theta \\ 0 \end{pmatrix} v
            + \begin{pmatrix} 0 \\ 0 \\ 1 \end{pmatrix} \omega.
\end{equation}

\paragraph{Controllability.}
Compute the Lie bracket of the two input vector fields with
\cref{eq:lie-bracket}:
\[
[g_1, g_2](q) = \frac{\partial g_2}{\partial q} g_1
              - \frac{\partial g_1}{\partial q} g_2
              = - \begin{pmatrix}
                    0 & 0 & -\sin\theta \\
                    0 & 0 & \cos\theta \\
                    0 & 0 & 0
                  \end{pmatrix}
                  \begin{pmatrix} 0 \\ 0 \\ 1 \end{pmatrix}
              = \begin{pmatrix} \sin\theta \\ -\cos\theta \\ 0 \end{pmatrix},
\]
since \( g_2 \) is constant. This is the unit vector orthogonal to the
sagittal axis, called the \textbf{slide vector field}: a velocity that
is immediately unrealizable, but that specifies a direction achievable
through a sequence of motions (a \emph{wriggle}, see
\cref{sec:lie-bracket-interpretation}). It is not a linear combination
of \( g_1 \) and \( g_2 \). Hence
the construction of the accessibility distribution ends immediately
with \( \Delta_{\mathcal{A}} = \Delta_2 \) and
\[
\dim \Delta_{\mathcal{A}} = \operatorname{rank}
\begin{pmatrix} g_1 & g_2 & [g_1, g_2] \end{pmatrix} = 3 = n.
\]
The unicycle is controllable with degree of nonholonomy
\( \kappa = 2 \), and constraint \cref{eq:rolling} is completely
nonholonomic, as already shown constructively with the universal
maneuver.

\begin{nb}[Robots Kinematically Equivalent to the Unicycle]
A unicycle in the strict sense is not statically balanced. Two
statically balanced structures are \emph{kinematically} equivalent to
it (their mechanics is different):
\begin{itemize}
    \item \textbf{Differential drive}: with wheel radii \( r \),
          angular velocities \( \omega_R, \omega_L \), and distance
          \( d \) between the contact points,
          \[
          v = r \frac{\omega_R + \omega_L}{2},
          \qquad
          \omega = (\omega_R - \omega_L) \frac{r}{d}.
          \]
          There is a one-to-one correspondence between \( (v, \omega) \)
          and \( (\omega_R, \omega_L) \). Here \( (x, y) \) is the
          midpoint between the contact points of the two fixed wheels
          and \( \theta \) is the orientation of the body.
    \item \textbf{Synchro-drive}: the two motors directly control
          \( v \) and \( \omega \), common to the three steerable
          wheels, where \( (x, y) \) can be any point of the robot
          (e.g.\ the geometric center of the three contact points).
          Unlike the differential drive, the orientation of the body
          never changes, unless a third actuator is added.
\end{itemize}
\end{nb}

\subsection{The Bicycle Model}
\label{subsec:bicycle-model}

The bicycle has one fixed rear wheel and one steerable front wheel.
Its configuration is \( q = (x, y, \theta, \phi)^\top \), with
\( (x, y) \) the position of the rear wheel contact point, \( \theta \)
the orientation of the body, and \( \phi \) the steering angle. The
configuration space is
\( C = \mathbb{R}^2 \times SO(2) \times SO(2) \), so \( n = 4 \).

Let \( l \) be the distance between the contact points of the two
wheels. The position of the front wheel is
\begin{align}
    x_f &= x + l \cos\theta, \\
    y_f &= y + l \sin\theta.
\end{align}
The two pure rolling constraints (rear and front wheel) are
\begin{align}
    \sin\theta \, \dot{x} - \cos\theta \, \dot{y} &= 0, \\
    \sin(\theta + \phi) \, \dot{x}_f
    - \cos(\theta + \phi) \, \dot{y}_f &= 0.
\end{align}
Each wheel has its own zero-motion line. The two zero-motion lines meet
at the \textbf{instantaneous center of rotation} (ICR), whose position
changes with the configuration; every point of the vehicle moves
instantaneously along an arc of circle centered at the ICR.

Substituting \( \dot{x}_f, \dot{y}_f \), the second constraint becomes
\[
\sin(\theta + \phi) \, \dot{x}
- \cos(\theta + \phi) \, \dot{y}
- l \cos\phi \, \dot{\theta} = 0.
\]
In matrix form:
\[
\underbrace{
\begin{pmatrix}
    \sin\theta & -\cos\theta & 0 & 0 \\
    \sin(\theta + \phi) & -\cos(\theta + \phi) & -l\cos\phi & 0
\end{pmatrix}}_{A(q)}
\begin{pmatrix} \dot{x} \\ \dot{y} \\ \dot{\theta} \\ \dot{\phi} \end{pmatrix} = 0.
\]
The matrix has constant rank \( k = 2 \), so the null space has
dimension \( m = 4 - 2 = 2 \). There are two choices of basis with a
physical meaning, depending on which wheel is driven.

\paragraph{Rear-wheel drive.}
Look for a first vector of the form
\( g_1 = (\cos\theta, \sin\theta, \alpha, 0)^\top \) (motion of the rear
wheel along its sagittal axis). Imposing the second constraint gives
\( \sin\phi - l\alpha\cos\phi = 0 \), i.e.\
\( \alpha = \frac{1}{l}\tan\phi \). With \( g_2 \) the pure steering
velocity, a basis is
\[
g_1^R(q) =
\begin{pmatrix}
    \cos\theta \\ \sin\theta \\ \frac{1}{l}\tan\phi \\ 0
\end{pmatrix},
\qquad
g_2^R(q) =
\begin{pmatrix}
    0 \\ 0 \\ 0 \\ 1
\end{pmatrix},
\]
and the kinematic model is
\begin{equation}
\label{eq:bicycle-rear}
    \dot{q} =
    \begin{pmatrix}
        \cos\theta \\ \sin\theta \\ \frac{1}{l}\tan\phi \\ 0
    \end{pmatrix} v
    + \begin{pmatrix} 0 \\ 0 \\ 0 \\ 1 \end{pmatrix} \omega.
\end{equation}
Here \( v = \pm\sqrt{\dot{x}^2 + \dot{y}^2} \) is the velocity of the
rear wheel contact point, and \( \omega = \dot{\phi} \) is the steering
velocity.

\paragraph{Singularity.}
The model \cref{eq:bicycle-rear} is singular at \( \phi = \pm\pi/2 \),
where \( \tan\phi \) is undefined and \( g_1^R \) is not defined: the
model assumes pure rolling, and at \( \phi = \pm\pi/2 \) any non-zero
\( v \) would make the front wheel slide. In practice, the singularity
is avoided by limiting the steering angle.

\paragraph{Controllability (rear-wheel drive).}
The first two brackets are
\[
g_3^R = [g_1^R, g_2^R] =
\begin{pmatrix} 0 \\ 0 \\ -\frac{1}{l\cos^2\phi} \\ 0 \end{pmatrix},
\qquad
g_4^R = [g_1^R, g_3^R] =
\begin{pmatrix}
    -\frac{\sin\theta}{l\cos^2\phi} \\
    \frac{\cos\theta}{l\cos^2\phi} \\ 0 \\ 0
\end{pmatrix}.
\]
They are obtained as follows. Since \( g_2^R = e_4 \) is constant,
\( g_3^R = -\partial g_1^R / \partial \phi \). The field \( g_3^R \)
depends only on \( \phi \), and the \( \phi \)-component of
\( g_1^R \) is zero, so \( (\partial g_3^R/\partial q)\, g_1^R = 0 \);
moreover \( g_3^R \) has only the \( \theta \)-component, therefore
\( g_4^R = -(\partial g_1^R/\partial\theta)\, g_3^R \). The field
\( g_3^R \) changes only the heading (a rotation in place of the
vehicle), while \( g_4^R \) is parallel to the slide vector field of
the unicycle (a lateral displacement of the rear wheel).

Since \( g_1^R, g_2^R, g_3^R, g_4^R \) are linearly independent, with
\[
\det \begin{pmatrix} g_1^R & g_2^R & g_3^R & g_4^R \end{pmatrix}
= \frac{1}{l^2\cos^4\phi} \neq 0 \qquad (\phi \neq \pm\pi/2),
\]
we have \( \Delta_{\mathcal{A}}^R = \Delta_3^R \) and
\( \dim \Delta_{\mathcal{A}}^R = 4 = n \). The rear-wheel-drive bicycle
is controllable with degree of nonholonomy \( \kappa = 3 \), and the
rolling constraints are completely nonholonomic. The rank is lost
exactly at the singularity of the model.

\paragraph{Front-wheel drive.}
If the robot is driven by the front wheel, replace \( g_1^R \) with
\( g_1^F = g_1^R \cos\phi \):
\begin{equation}
\label{eq:bicycle-front}
\dot{q} =
\underbrace{
\begin{pmatrix}
    \cos\theta \cos\phi \\ \sin\theta \cos\phi \\ \frac{1}{l}\sin\phi \\ 0
\end{pmatrix}}_{g_1^F} v
+ \underbrace{\begin{pmatrix} 0 \\ 0 \\ 0 \\ 1 \end{pmatrix}}_{g_2^F} \omega.
\end{equation}
Now \( v = \pm\sqrt{\dot{x}_f^2 + \dot{y}_f^2} \) is the velocity of the
front wheel contact point, and the model is \emph{not} singular at
\( \phi = \pm\pi/2 \): in that situation a front-wheel-drive bicycle
simply pivots around the rear wheel. The brackets are
\[
g_3^F = [g_1^F, g_2^F] =
\begin{pmatrix}
    \cos\theta \sin\phi \\ \sin\theta \sin\phi \\ -\frac{\cos\phi}{l} \\ 0
\end{pmatrix},
\qquad
g_4^F = [g_1^F, g_3^F] =
\begin{pmatrix}
    -\frac{\sin\theta}{l} \\ \frac{\cos\theta}{l} \\ 0 \\ 0
\end{pmatrix}.
\]
Also \( g_1^F, g_2^F, g_3^F, g_4^F \) are linearly independent, so the
front-wheel-drive bicycle is controllable with \( \kappa = 3 \).

\begin{nb}[Instantaneous Center of Rotation and Ackermann Steering]
The two zero-motion lines (the directions in which each wheel cannot
move) intersect at the \textbf{instantaneous center of rotation
(ICR)}, the point around which the robot is rotating at a given
instant. For a tricycle, the ICR is unique. For a car-like robot with
two steerable front wheels, the two front zero-motion lines may
intersect the rear zero-motion line at different points, so a single
ICR does not exist. The \textbf{Ackermann steering} mechanism solves
this by automatically orienting the front wheels so that their
zero-motion lines intersect at a single ICR.
\end{nb}

\begin{nb}[Kinematically Equivalent to the Bicycle]
The bicycle is not statically balanced (its support polygon is a line
segment). The \textbf{tricycle} is modeled as a bicycle by placing
\( (x, y) \) at the midpoint between the two rear wheels, with \( \theta \)
the orientation of the vehicle and \( \phi \) the steering angle. The
\textbf{car-like} robot is modeled as a bicycle once Ackermann steering
is in place. In both cases the kinematic model is
\cref{eq:bicycle-rear} or \cref{eq:bicycle-front}, depending on whether
the drive is on the rear or on the front wheels.
\end{nb}

\section{Chained Form}
\label{sec:chained-form}

Many kinematic models of mobile robots can be recast in the same
canonical form. This is useful because it allows the use of systematic
procedures for planning and control. We limit the analysis to systems
with two inputs, such as the unicycle and the bicycle.

\begin{definition}[Chained Form]
A \( (2, n) \) \textbf{chained form} is the driftless system
\begin{equation}
\label{eq:chained}
\begin{aligned}
    \dot{z}_1 &= w_1, \\
    \dot{z}_2 &= w_2, \\
    \dot{z}_3 &= z_2 w_1, \\
              &\;\;\vdots \\
    \dot{z}_n &= z_{n-1} w_1,
\end{aligned}
\end{equation}
with \( z = (z_1, \dots, z_n) \) the \( n \)-dimensional state and
\( w = (w_1, w_2) \) the 2-dimensional input.
\end{definition}

\begin{proposition}
\label{prop:chained-controllable}
The \( (2,n) \) chained form is controllable, with degree of
nonholonomy \( \kappa = n - 1 \).
\end{proposition}

\begin{proof}
The input fields are \( g_1 = (1, 0, z_2, \dots, z_{n-1})^\top \) and
\( g_2 = e_2 \). In the Jacobian of \( g_1 \), the entry on row
\( i \ge 3 \) depends on \( z_{i-1} \), so the column of \( z_j \)
(\( 2 \le j \le n-1 \)) is \( e_{j+1} \). Since \( g_2 \) is constant,
\[
[g_1, e_j] = -\frac{\partial g_1}{\partial z}\, e_j = -e_{j+1},
\qquad 2 \le j \le n-1.
\]
Starting from \( g_2 = e_2 \), the iterated brackets are
\[
\text{ad}_{g_1}^k g_2 = (-1)^k e_{k+2}, \qquad k = 0, \dots, n-2 .
\]
Together with \( g_1 \), which has first component 1, they span
\( \mathbb{R}^n \). Each bracket adds exactly one dimension, so the
rank \( n \) is reached with brackets of length \( n - 1 \) and not
before: \( \kappa = n - 1 \).
\end{proof}

If \( w \) is constant, the chained form is a \emph{linear} system (for
\( w_1 \) constant, each equation \( \dot z_i = z_{i-1} w_1 \) is linear
in \( z \)), so in digital control, where inputs are piecewise
constant, the system is piecewise linear.

A generic \( n \)-dimensional two-input driftless system
\begin{equation}
\label{eq:two-input}
    \dot{q} = G(q) u = g_1(q) u_1 + g_2(q) u_2
\end{equation}
can be transformed into the chained form \cref{eq:chained} via a
coordinate transformation and an input transformation,
\[
z = \psi(q), \qquad u = \Gamma(q)\, w,
\]
under necessary and sufficient conditions (not treated here). In
particular, systems of the form \cref{eq:two-input} with \( n \le 4 \)
can always be put in chained form. This applies to the unicycle and the
bicycle, but the technique does not extend automatically to larger
systems.

\subsection{The Unicycle in Chained Form}
\label{subsec:unicycle-chained}

Consider the unicycle model \cref{eq:unicycle}. Use the coordinate
transformation
\begin{equation}
\label{eq:chained-coord-transform}
z_1 = \theta, \qquad
z_2 = x\cos\theta + y\sin\theta, \qquad
z_3 = x\sin\theta - y\cos\theta,
\end{equation}
and the input transformation
\[
v = z_3 w_1 + w_2, \qquad \omega = w_1.
\]
The model becomes the \( (2,3) \) chained form
\begin{align}
    \dot{z}_1 &= w_1, \\
    \dot{z}_2 &= w_2, \\
    \dot{z}_3 &= z_2 w_1.
\end{align}

\paragraph{Verification.}
First, \( \dot z_1 = \dot\theta = \omega = w_1 \). Then
\begin{align}
    \dot{z}_2
    &= \dot{x} \cos\theta + \dot{y} \sin\theta
       - x \sin\theta\, \dot{\theta} + y \cos\theta\, \dot{\theta} \\
    &= v - \underbrace{\dot{\theta}}_{\omega}
       \underbrace{(x \sin\theta - y \cos\theta)}_{z_3}
     = v - \omega z_3 = w_2 ,
\end{align}
and
\begin{align}
    \dot{z}_3
    &= \dot{x} \sin\theta - \dot{y} \cos\theta
       + x \cos\theta\, \dot{\theta} + y \sin\theta\, \dot{\theta} \\
    &= \underbrace{(\dot{x} \sin\theta - \dot{y} \cos\theta)}_{0
       \text{ by the constraint}}
       + \dot{\theta} (x \cos\theta + y \sin\theta)
     = \omega\, z_2 = z_2 w_1 .
\end{align}

Note that \( z_1 \) is simply the orientation \( \theta \), while
\( z_2 \) and \( z_3 \) are the position of the unicycle in a frame
whose first axis is the sagittal axis of the vehicle,
\( (\cos\theta, \sin\theta) \), and whose second axis is the
right-hand normal \( (\sin\theta, -\cos\theta) \)
(\cref{fig:z2-z3}).

\begin{figure}[ht]
    \centering
    \includegraphics[width=0.5\linewidth]{images/z2z3.eps}
    \caption{Interpretation of \( z_2 \) and \( z_3 \): coordinates of
    the unicycle in the frame aligned with the sagittal axis (first
    axis) and its right-hand normal (second axis).}
    \label{fig:z2-z3}
\end{figure}

\subsection{The Bicycle in Chained Form}
\label{subsec:bicycle-chained}

The rear-wheel-drive bicycle model \cref{eq:bicycle-rear} can also be
transformed into chained form. Use the coordinate transformation
\[
z_1 = x, \quad
z_2 = \frac{1}{l}\sec^3\theta \tan\phi, \quad
z_3 = \tan\theta, \quad
z_4 = y,
\]
and the input transformation
\[
v = \frac{w_1}{\cos\theta}, \qquad
\omega = -\frac{3}{l}\, w_1 \sec\theta \tan\theta \sin^2\phi
        + l\, w_2 \cos^3\theta \cos^2\phi.
\]
The resulting system is the \( (2,4) \) chained form
\begin{align}
    \dot{z}_1 &= w_1, \\
    \dot{z}_2 &= w_2, \\
    \dot{z}_3 &= z_2 w_1, \\
    \dot{z}_4 &= z_3 w_1.
\end{align}
The input transformation is undefined at \( \cos\theta = 0 \), i.e.\
at \( \theta = \pi/2 + k\pi \): the equivalence between the two
models is subject to \( \theta \neq \pi/2 + k\pi \).

\paragraph{Verification.}
With \( \dot\theta = v\tan\phi/l = w_1\tan\phi/(l\cos\theta) \) and
\( \dot\phi = \omega \), the equations that do not involve \( \omega \)
are checked directly:
\begin{align}
    \dot z_1 &= v\cos\theta = w_1, \\
    \dot z_3 &= \sec^2\theta\, \dot\theta
              = \frac{1}{l}\sec^3\theta\tan\phi\, w_1 = z_2 w_1, \\
    \dot z_4 &= v\sin\theta = w_1\tan\theta = z_3 w_1 .
\end{align}
The input \( \omega \) enters only through \( \dot z_2 \):
\[
\dot z_2
= \frac{3}{l^2}\sec^4\theta\tan\theta\tan^2\phi\; w_1
  + \frac{1}{l}\sec^3\theta\sec^2\phi\;\omega .
\]
Imposing \( \dot z_2 = w_2 \) and solving for \( \omega \), with
\( \cos^3\theta\sec^4\theta = \sec\theta \) and
\( \cos^2\phi\tan^2\phi = \sin^2\phi \), gives the input
transformation above.

\begin{nb}[Trailed Vehicles]
An interesting property of \textbf{trailed vehicles} emerges when
considering a car-like robot: if the trailer is hitched exactly at the
midpoint between the two rear wheels, the kinematic model can still be
converted into chained form. This property holds recursively: if the
trailer itself has two rear wheels and an additional trailer is
hitched at their midpoint, the overall kinematic model can still be
expressed in chained form.
\end{nb}

\section{Takeaways}
\label{sec:wmr-mechanics-takeaways}

\begin{enumerate}
    \item A wheeled robot is in \textbf{static balance} if the
          projection of its center of mass falls inside the
          \textbf{support polygon}, the convex hull of its contact
          points. Unicycle and bicycle are not balanced; differential
          drive, synchro-drive, tricycle and car-like robots are.
    \item Three types of wheels exist: \textbf{fixed}, \textbf{steerable},
          and \textbf{caster}. They differ in whether they can steer,
          whether they are actuated, and whether they are kinematically
          ``invisible''.
    \item The main kinematic structures are the \textbf{differential
          drive}, the \textbf{synchro-drive}, the \textbf{tricycle}, and
          the \textbf{car-like} robot.
    \item The \textbf{mechanical differential} allows two wheels on the
          same axle to rotate at different speeds, avoiding slip during
          turns.
    \item Geometric constraints \( h(q) = 0 \) reduce the configuration
          space to a sub-manifold; kinematic constraints in Pfaffian
          form \( A(q)\dot q = 0 \) restrict the admissible velocities
          to \( \mathcal{N}(A(q)) \). Integrable kinematic constraints
          are holonomic, the others are nonholonomic.
    \item The kinematic model \( \dot q = G(q) u \) is driftless, and
          \( u \) are pseudo-velocities. The constraints are completely
          nonholonomic if and only if the accessibility rank condition
          \( \dim \Delta_{\mathcal{A}} = n \) holds, i.e.\ if the model
          is controllable.
    \item The unicycle (\( \kappa = 2 \)) and the bicycle
          (\( \kappa = 3 \), singular at \( \phi = \pm\pi/2 \) only in
          the rear-wheel-drive case) are controllable; differential
          drive and synchro-drive are kinematically equivalent to the
          unicycle, tricycle and car-like to the bicycle.
    \item Both can be put in \textbf{chained form}, always controllable
          with \( \kappa = n - 1 \), which allows systematic planning
          and control.
\end{enumerate}